\documentclass[main]{subfiles}

\title[AutoML: Dynamic Configuratino]{Chapter 11 -- Quiz}
\subtitle{Multiple Choice Questions}
\author[André Biedenkapp]{André Biedenkapp}
\date{}

\titlebackground*{images/AutoML_LearningToLearn2}

\newcommand{\correct}[1]{\textcolor{ForestGreen}{\textbf{#1} \checkmark}}
\newcommand{\wrong}[1]{#1}

\begin{document}

    \maketitle

%----------------------------------------------------------------------

%----------------------------------------------------------------------
\begin{frame}{Q0 \hfill \small\normalfont\textit{(single correct answer)}}

\textbf{DAC is a black-box optimization approach.}

\begin{itemize}
    \item \wrong{(A) \; True.}
    \item \correct{(B) \; False.}
\end{itemize}

\end{frame}
%-------------------------------------------------
%----------------------------------------------------------------------
\begin{frame}{Q1 \hfill \small\normalfont\textit{(single correct answer)}}

\textbf{Static configuration is guaranteed to outperform DAC for small budgets.}

\begin{itemize}
    \item \wrong{(A) \; True.}
    \item \correct{(B) \; False.}
\end{itemize}

\end{frame}
%----------------------------------------------------------------------

%----------------------------------------------------------------------
\begin{frame}{Q2 \hfill \small\normalfont\textit{(single correct answers)}}

\textbf{DAC can be solved by optimization}

\begin{itemize}
    \item \correct{(A) \; True.}
    \item \wrong{(B) \; False.}
\end{itemize}

\end{frame}
%----------------------------------------------------------------------

%----------------------------------------------------------------------
\begin{frame}{Q3 \hfill \small\normalfont\textit{(single correct answer)}}

\textbf{DAC can \textbf{not} be solved with RL.}

\begin{itemize}
    \item \wrong{(A) \; True.}
    \item \correct{(B) \; False.}
\end{itemize}

\end{frame}
%----------------------------------------------------------------------

%----------------------------------------------------------------------
\begin{frame}{Q4 \hfill \small\normalfont\textit{(single correct answers)}}

\textbf{Learning curve extrapolation for HPO is a DAC problem}

\begin{itemize}
    \item \wrong{(A) \; True.}
    \item \correct{(B) \; False.}
\end{itemize}

\end{frame}
%----------------------------------------------------------------------

%----------------------------------------------------------------------
\begin{frame}{Q5 \hfill \small\normalfont\textit{(multiple correct answers)}}

\textbf{Iterative optimization heuristics ... \textit{(Select all that apply)}}

\begin{itemize}
    \item \wrong{(A) \; use the black-box view.}
    \item \correct{(B) \; propose a (set of) solution candiate(s) in each iteration.}
    \item \correct{(C) \; aim at improving the solution quality in each iteration.}
    \item \wrong{(D) \; only take the final solution quality into account.}
\end{itemize}

\end{frame}
%----------------------------------------------------------------------

%----------------------------------------------------------------------
\begin{frame}{Q6 \hfill \small\normalfont\textit{(multiple correct answer)}}

\textbf{Population-based training (PBT) ... \textit{(Select all that apply)}}

\begin{itemize}
    \item \wrong{(A) \; changes only a single hyperparamter in each iteration.}
    \item \correct{(B) \; drops the worst performing configurations in each iteration.}
    \item \correct{(C) \; inherits weights from well performing networks.}
    \item \wrong{(D) \; uses mutation to improve the performance of poorly performing configs.}
\end{itemize}

\end{frame}
%----------------------------------------------------------------------

%----------------------------------------------------------------------
\begin{frame}{Q7 \hfill \small\normalfont\textit{(multiple correct answers)}}

\textbf{Which PBT variants are capable of jointly tuning hyperparameters and architectures? ... \textit{(Select all that apply)}}

\begin{itemize}
    \item \wrong{(A) \; PBT.}
    \item \wrong{(B) \; Population-based Bandits (PB2).}
    \item \correct{(C) \; Bayesian-Generational PBT.}
    \item \wrong{(D) \; Meta-PB2.}
\end{itemize}

\end{frame}
%----------------------------------------------------------------------

%----------------------------------------------------------------------
\begin{frame}{Q8 \hfill \small\normalfont\textit{(multiple correct answers)}}

\textbf{To combine Bayesian Optimization (BO) and Population-based Training (PBT), we could ... \textit{(Select all that apply)}}

\begin{itemize}
    \item \correct{(A) \; use an asynchronous parallel BO.}
    \item \correct{(B) \; use BO with a nonstationary model.}
    \item \wrong{(C) \; predict how well a configuration performed in the past.}
    \item \correct{(D) \; try to predict how well a configuration will do in the future.}
\end{itemize}

\end{frame}
%----------------------------------------------------------------------

%----------------------------------------------------------------------
\begin{frame}{Q9 \hfill \small\normalfont\textit{(multiple correct answers)}}

\textbf{Bayesian generational PBT ... \textit{(Select all that apply)}}

\begin{itemize}
    \item \correct{(A) \; samples novel architectures in every iteration.}
    \item \wrong{(B) \; copies weights across dissimilar architectures and makes use of isomorphisms.}
    \item \correct{(C) \; transfers knowledge across generations via distillation.}
    \item \wrong{(D) \; purely cares about architectural hyperparameters.}
\end{itemize}

\end{frame}
%----------------------------------------------------------------------

%----------------------------------------------------------------------
\begin{frame}{Q10 \hfill \small\normalfont\textit{(multiple correct answers)}}

\textbf{Dynamic Algorithm Configuration ... \textit{(Select all that apply)}}

\begin{itemize}
    \item \correct{(A) \; can theoretically give exponential improvements over static HPO.}
    \item \correct{(B) \; is too expensive for many real-world applications.}
    \item \correct{(C) \; can produce policies that can solve heterogeneous problem instances.}
    \item \wrong{(D) \; should only ever be used for single hyperparameters.}
\end{itemize}

\end{frame}
%----------------------------------------------------------------------

%----------------------------------------------------------------------
\begin{frame}{Q11 \hfill \small\normalfont\textit{(single correct answers)}}

\textbf{Is it possible to learn dynamic configuration schedules online, i.e. not in a dedicated offline learning phase?}

\begin{itemize}
    \item \correct{(A) \; True.}
    \item \wrong{(B) \; False.}
\end{itemize}

\end{frame}
%----------------------------------------------------------------------

%----------------------------------------------------------------------
\begin{frame}{Q12 \hfill \small\normalfont\textit{(single correct answers)}}

\textbf{RL can learn general policies only if the policy is explicitly aware of the environment context?}

\begin{itemize}
    \item \wrong{(A) \; True.}
    \item \correct{(B) \; False.}
\end{itemize}

\end{frame}
%----------------------------------------------------------------------

%----------------------------------------------------------------------
\begin{frame}{Q13 \hfill \small\normalfont\textit{(multiple correct answers)}}

\textbf{Dynamic configuration can be achieved via ... \textit{(Select all that apply)}}

\begin{itemize}
    \item \correct{(A) \; Heuristics.}
    \item \correct{(B) \; Black-Box Optimization.}
    \item \correct{(C) \; Evolution.}
    \item \correct{(D) \; RL.}
\end{itemize}

\end{frame}

%----------------------------------------------------------------------
\begin{frame}{Q14 \hfill \small\normalfont\textit{(multiple correct answers)}}

\textbf{Informative context features ... \textit{(Select all that apply)}}

\begin{itemize}
    \item \correct{(A) \; are often domain and problem specific.}
    \item \wrong{(B) \; are computed for every configuration step.}
    \item \correct{(C) \; characterize the problem.}
    \item \wrong{(D) \; change during the run of a target algorithm.}
\end{itemize}

\end{frame}
%----------------------------------------------------------------------

\end{document}